\documentclass[lang=cn,10pt]{elegantbook}
\usepackage{hyperref}
\usepackage{tikz}
\usetikzlibrary{arrows.meta, positioning}
\usepackage{pgfplots}
\pgfplotsset{compat=1.18}
\usepgfplotslibrary{statistics}
\usepackage{float}

% 让每个 \chapter 都从奇数页（右侧页）开始
\makeatletter
\@twosidetrue
\@mparswitchtrue
\typeout{DEBUG twoside = \if@twoside YES\else NO\fi}
\makeatother



\title{Statistic Beans：統計学豆知識}
\subtitle{不应错过的统计学基础知识}

\author{郑在时}
% \institute{Elegant\LaTeX{} Program}
\date{ 2026/3/4}
\version{0.3ed   {Typeset with \LaTeX{}}}
% \bioinfo{自定义题头}{信息}

\extrainfo{不要以为抹消过去，重新来过，即可发生什么改变。—— 比企谷八幡}

\setcounter{tocdepth}{3}

% \logo{logo-blue.png}
\cover{cover.jpg}

% 本文档命令
\usepackage{array}
\newcommand{\ccr}[1]{\makecell{{\color{#1}\rule{1cm}{1cm}}}}

% 修改标题页的橙色带
\definecolor{customcolor}{RGB}{255,255,120}
\colorlet{coverlinecolor}{customcolor}

\begin{document}

\maketitle

\frontmatter


\newpage
\thispagestyle{empty}
\vspace*{\fill}
\begin{center}
    \large \textbf{\textit{To my friends.}}
\end{center}
\vfill  


\tableofcontents
\mainmatter



\chapter{写在前面}

本知识手册记录了一些容易被忽视的统计基础概念，这也是为什么取名为“Statistic Beans：統計学豆知識”。如果您在阅读时遇到了困难，欢迎联系作者。

作者的联系方式：
\begin{itemize}
  \item 邮箱：\email{zzs234@yeah.net}
  \item Blog：\href{https://zes4u.pages.dev}{https://zesforu.pages.dev}
  \item GitHub：\href{https://github.com/Zes-Z}{https://github.com/Zes-Z}
\end{itemize}

\section{本书内容}
\begin{itemize}
  \item 第1章，也就是你正在阅读的这里，会简要介绍2-8章的关键内容;
  \item 第2章，描述了什么是统计，\textbf{统计}到底在干什么?
  \item 第3章，一个完整的统计过程!
  \item 第4章，主要回答：如何将统计得到的\textbf{一堆数字}优雅地\textbf{呈现}出来？并最大效率地\textbf{描述}这堆数字表达的\textbf{信息}？
  \item 第5章，更进一步，我们需要一些\textbf{概率}和\textbf{分布}；
  \item 第6章，假设检验：\textbf{从数字}形式的统计量\textbf{到文字}描述!\\
  如果我们假设今天喜家德没有喜三鲜，结果刚坐下点单时，发现隔壁桌吃的都是喜三鲜!\\
  啊哈，看来我们的假设只有非常小的概率成立(因为可能有人外带了喜三鲜！)\\
  所以，我们有很大把握今天喜家德有喜三鲜;\\
  顺带一提，其实西芹鲜肉也不错!
  \item 第7章，方差分析，分析的是三个以上的\textbf{组均值是否存在显著差异}，差异的\textbf{原因}是什么；
  \item 第8章，\textbf{给出}多组一个或多个变量的\textbf{对应数值}，能否预测我们\textbf{想知道的数据}？\\
  例如：给你最近十天的\textbf{日照时间--平均气温}的对应数据，你能根据今天的日照时间预测今天的平均气温吗?\\
  不可以看天气预报!
  \item 第9章，
  \item 附录部分放置了一些常见分布的概率临界值表供参考
\end{itemize}


\section{记录契机}
这本“知识手册”，总结了一些学习概率论与数理统计的心得，足以覆盖初入统计涉及到的一些基础概念。其实在本学期开学前，作者本以为能从从容容地应付下
述课程，结果却让我大为惊叹。此时距离《语言研究中的统计学》这门课结束还有半个月，距离托
业考试还有一月余五天，最近无心学习，所以才有时间写这本小册子。第一次尝试写知识手册，不过请读者诸位放心，全文无杜撰，无毒，
大可放心“食用”。

那么，为什么写这本小册子？答案很简单：为了方便。就如同日语中为了发音便利
而产生了「音便」一般。\textbf{如果把学习过的知识，总结记录下来，不论是对于再次查阅还是加深理解，都很有益。}数周前，一位朋友和我探讨一些统计上的问题，大概是关于$F$分布和$p$
值之类的，我本以为能够清晰简明地解释，但并非如此；几天前，和另一位
朋友闲聊时，涉及到“统计学”时，不禁感叹道：“本来**就没啥说服力,全靠吹气球”，如果不引
用一点方法，那便真正陷入了循环论证。于是我再次翻开了参考书籍，萌发了将部分所学记录下来的念头。

此书写给自己和我的朋友们，如果有些许助益的话，也算是满足了我小小的幸福感。


\section{阅读建议}
以下是一些阅读本书的建议：
\begin{itemize}
  \item 出现数学公式时，不会作过多的术语解说，只需要知道这串符号在\textbf{算些什么}，读者也可查阅相关资料。
  \item 本书会在术语首次出现时解释说明，后不再赘述，建议添加标签、笔记。
  \item 符号标记为分析案例，建议动手尝试，不要只看不做，那没有用。
  \item 由于作者水平、认知有限，特殊标记的内容仅为作者的见解，建议读者自查相关资料。
  \item 感到迷惑时立即停下，你需要补充相关知识，或者回头看看。基础不牢，地动山摇。
  \item 不想读时立即停下，你需要休息。
\end{itemize}






\chapter{数据与统计量}
\section{数据一定是数字吗?}
有人说，我们身处于一个物质世界。

但其实，我们也生活在一个数据世界之中...我们的大脑，无时不刻都在处理海量的信息。每天早晨在睡意朦胧中，你听到闹钟响，然后爬起身来；洗漱时，你偶尔也会发呆，或回忆昨天发生的事情；打开手机和电脑，你快速地浏览消息并准备回复；工作时，你盯着满屏的文字、图片、excel表格，感到两眼发花...

相信你已经感受到，\textbf{数据未必以数字的形式出现}。其实它也可以是文字、符号、甚至声音、图片、视频等，不过为了方便处理和研究，数据最终都会被转化为数字或类数字形式。
\begin{definition}[数据] 
数据是对研究对象的某种属性或特征进行观测，测量或记录后得到的、能够用于描述和分析的信息。

数据是随机现象中观测结果的集合，是统计分析的基本材料。
\end{definition}



\section{统计量是什么?}
获得足够量的数据以后，便有了进行统计研究的基础支撑。但是，我们不能简单地罗列出所有数据，这非常的不便，也没有意义，当然这也不能被称为\textbf{统计}。

因此，我们需要站在更高的总体的视角去提取样本数据(集)的信息，进行描述，分析；同时也应该以更加细致入微的视角，观察每个样本数据(集)之间的差异、联系。
\begin{definition}[统计量] 
统计量是由样本数据计算得到的、且不包含未知总体参数的函数。

统计量 = 对样本数据进行某种计算得到的数值，用来描述样本或推断总体。
\end{definition}

常见的统计量用于基础描述统计，主要有:
\begin{itemize}
  \item \textbf{样本均值}$\bar{x}$：把所有数据加起来再除以个数，反映数据的“平均水平”或中心位置。例如，5名同学的日语成绩分别是 60、70、80、90、100 分，则 $\bar{x}=(60+70+80+90+100)/5=80$ 分。
  \item \textbf{中位数}：把数据从小到大排列后位于正中间的数。它比均值更“扛得住”极端值。例如上面的成绩排序后中间是 80，中位数就是 80 分。
  \item \textbf{众数}：一组数据中出现次数（频数）最多的取值，代表“最常见的取值”。例如成绩 60、70、70、80、90 中，70 出现了两次、次数最多，所以众数是 70。一组数据可能没有众数，也可能有多个众数（如出现次数并列最多时）。
  \item \textbf{极差}：最大值减去最小值，是最简单的波动范围度量。例如最高分 100、最低分 60，极差就是 $100-60=40$ 分。它计算简单，但很容易被极端值带偏。
  \item \textbf{方差}：各数据与均值之差的平方的平均值，衡量数据整体偏离均值的程度。方差越大，说明数据越“散”。
  \item \textbf{标准差}：方差的算术平方根，和原始数据单位一致，所以比方差更直观。它衡量数据围绕均值的离散程度：标准差越大，数据越分散；标准差越小，数据越集中。例如成绩的标准差是 10 分，可以理解为“成绩围绕平均分波动的典型幅度约为 10 分”（对近似正态分布的数据，约 68\% 落在均值$\pm 1$个标准差之内）。
  \item \textbf{分位数}：把排序后的数据按位置切成若干等份。中位数就是第 50\% 分位数（0.5 分位数）；第 25\%、第 75\% 分位数分别称为下四分位数和上四分位数，用来描述数据分布在哪些位置。
  \item \textbf{相关系数}：衡量两个变量之间线性相关的方向和强弱，取值范围在 $[-1,1]$。越接近 1 表示正相关越强，越接近 $-1$ 表示负相关越强，接近 0 表示几乎没有线性关系。
\end{itemize}


进阶的统计量用于从样本推断总体，主要有:
\begin{itemize}
  \item \textbf{标准化统计量}：把数据减去均值再除以标准差，得到标准分（$z$ 分数）。它把不同尺度的数据放到同一把“尺子”上，方便比较。例如小明日语考 85 分、全班均值 80、标准差 5，则 $z=(85-80)/5=1$，说明他比平均高 1 个标准差。
  \item \textbf{偏度}：衡量分布左右不对称的程度和方向。偏度大于 0 说明右尾更长（右偏），小于 0 说明左尾更长（左偏），等于 0 说明大致对称。
  \item \textbf{峰度}：衡量分布是“尖”还是“平”，反映数据在均值附近集中的程度以及尾部的厚薄。峰度越大，分布通常越“尖”、尾部越厚。
  \item \textbf{$t$统计量}：在总体方差未知、样本量较小时，用于检验均值是否等于某个值或两组均值是否有差异。它由“样本均值与假设值的差”除以“标准误差”得到，是 $t$ 检验的核心。
  \item \textbf{$F$统计量}：用于比较两个方差是否相等，或检验多个组的均值是否存在显著差异（方差分析）。它通常是两个方差之比，是 $F$ 检验和方差分析的核心。
  \item \textbf{$\chi^2$统计量}：用于检验分类数据的“实际频数”与“理论频数”是否吻合，例如检验某枚骰子是否公平、两个分类变量是否独立。它把每个类别的“实际值减期望值”的平方除以期望值后加总得到。
\end{itemize}



\section{统计 = 数据 + 统计量 + 统计推断}
如果说\textbf{数据}是统计的原材料，\textbf{统计量}是从大量的数据中提取，精简信息，
那么我们还需要一个\textbf{统计推断}，执行一个\textbf{由样本推测总体}的操作，
才能完成一个完整的统计，\textbf{统计推断}也是统计学最核心的部分。

其中，数据回答：“我们观察到了什么？”
统计量回答：“这些数据有什么特征？”
统计推断回答：“这些特征是否具有说服力，能否推广到更大的范围？”

用一个例子把三者串起来：假设你想了解全校学生的平均身高。
\begin{itemize}
  \item \textbf{数据}：随机抽取 100 名学生，记录下每个人的身高，这 100 个数字就是数据。
  \item \textbf{统计量}：计算这 100 个身高的平均值 $\bar{x}=168$ cm，这就是从数据中提取出的统计量。
  \item \textbf{统计推断}：用样本均值 $168$ cm 去估计全校学生的平均身高，并判断这个估计有多可靠、误差有多大。
\end{itemize}

\begin{remark}
  统计量是从样本中算出来的“已知数”，而总体的参数（如全校真实平均身高）是“未知数”。
  统计推断要做的事，就是借助统计量这个“已知数”，去合理地猜测那个“未知数”。
\end{remark}






\chapter{完整的统计过程}
\section{提出问题}
通常情况下，我们希望量化研究或说明某个问题时，才会使用到统计。我们收集的数据，需要计算的统计量，以及统计推断和描述都将紧紧围绕最初的问题展开，这时，作为统计过程起点的\textbf{问题提出}就显得尤为重要。

但也有时，我们需要先抛出一个大致的猜想，根据已有可行的统计方法、预计能收集到的数据，来反向不断明确我们的问题。这样就不会出现：你提出了一个天马行空，富有研究意义的课题，却不知从何下手，用统计的手段来验证和解释你的课题的情况。


\section{统计设计}
拿到一个研究问题后，先别急着冲出去收集数据。\textbf{统计设计}就是在动手之前，把“要什么数据、从哪里来、怎么来、来了以后怎么办”统统想清楚。设计做得好，后面的路才走得顺；设计做得糙，很可能收集完数据才发现方法不对、样本没代表性，白忙一场。

统计设计至少要想清楚以下几件事：
\begin{itemize}
  \item \textbf{研究问题与假设}：到底想回答什么？有没有一个可以检验的猜想？
  \item \textbf{总体与抽样对象}：我们的研究对象是谁？哪些人/物才算数？
  \item \textbf{抽样方法}：用概率抽样还是非概率抽样？怎么抽才能代表总体？
  \item \textbf{样本量}：大概需要多少条数据才够？太少了没说服力，太多了浪费。
  \item \textbf{变量与记录方式}：要记录哪些变量？怎么编码？用表格还是问卷？
  \item \textbf{统计检验的选择}：根据研究问题的类型判断该用哪种检验，并预判数据到手后会用哪些统计量。比如比较两组均值用$t$检验，比较多组均值用方差分析（$F$检验），考察分类变量是否独立用$\chi^2$检验，考察两个变量的关系用相关或回归。
\end{itemize}

例如：假设你是一位日语教师，想了解学校里日语学生的日语考试成绩。统计设计阶段就应该想清楚：总体是“全校日语学生”，抽样采用分层抽样（按年级、班级各抽一部分），记录学号、年级、成绩等变量；由于研究问题是“平均分是否达标”，属于比较均值，所以计划用均值、标准差来描述，再用$t$检验来判断。这样到了采集数据时，每一步都有章可循。

\begin{remark}
统计设计不是一次性的“填空题”。随着对问题的理解加深，设计常常需要回头调整——比如发现某个变量不好收集，或样本量预算不够，那就及时修改方案，总比硬着头皮收一堆用不上的数据强。
\end{remark}





\section{样本数据采集－－抽样方式与最佳组数}
在完成统计设计后，即可开始统计研究的第一步实际操作－－样本数据采集。

为了尽可能地保证统计的无偏差性和准确性，在采集样本数据时需要进行系统的、公平的、随机的抽样。既不能随心所欲地收集数据，也不能贪图方便省事。

例如：在上一个例子中，乍一看样本数据采集十分简单，只需收集一定量的日语成绩即可；但实际上
这里藏了不少坑！首先你不能仅收集自己班级学生的信息作为样本数据集，数据的来源要尽可能的全面，多样，无偏好，最终保证你收集到的那一小部分数据能够很好地代表总体。
所以你应该尽可能多的，随机收集学校里日语学生的日语考试成绩，再进行下一步的统计。

\begin{remark}
如果能收集到本学校所有日语学生的日语考试成绩自然是最好，这便是\textbf{总体}，但有时迫于现实原因或条件限制，我们只能收集到一部分数据，这便是用于\textbf{代表总体}的\textbf{样本}。
只要你从总体中抽取样本的方式足够科学，那么即使只收集到一小部分数据，也能够很好地反应总体的情况。
\end{remark}

以下是一些样本数据采集中使用到的经典抽样方法，前三种为概率抽样，后三种为非概率抽样：
\begin{itemize}
  \item 简单随机抽样
  \item 系统抽样/等距抽样
  \item 分层抽样
  \item 便利抽样
  \item 判断抽样
  \item 滚雪球抽样
\end{itemize}

若采集到$N$条数据，现要分组，以直方图(或其它)形式分组展示，则应将其分为$k$组，总的样本数$N$与组数$k$满足\textbf{斯特吉斯规则}。

\begin{theorem}[斯特吉斯规则 Sturges' rule]  
    \begin{center}
      $k = 1 + \log_2 N$
    \end{center}   

\end{theorem}


\section{数据的整理与描述}
初步采集的原始数据，往往是一堆乱糟糟的数字和记录，直接进行分析是不行的。\textbf{数据的整理与描述}，就是先把数据“收拾干净”，再用表格、图表和统计量把数据里的信息表现出来。

整理阶段通常要做这几件事：
\begin{itemize}
  \item \textbf{检查与清洗}：看看有没有缺失值、重复记录、明显录错的数（比如成绩写成 150 分、身高写成 1.8 米却录成 18 米）。
  \item \textbf{编码与分类}：把文字信息转成可计算的类别或数字，例如把“男/女”编码成 0/1。
  \item \textbf{排序与分组}：把数据按大小排列；需要画直方图时，按第 3.3 节的斯特吉斯规则确定组数再分组。
  \item \textbf{制作频数表}：统计每个取值或每个组里有多少条数据，一眼看出分布的大致形状。
  \item \textbf{可视化}：用直方图、箱型图、条形图等把分布展示出来（详见第 4 章）。
  \item \textbf{计算描述统计量}：均值、中位数、众数、方差、标准差、分位数等，用数字概括数据特征。
\end{itemize}

例如：你收集到 50 名日语学生的成绩，先检查有没有缺考或录错的分数；再按 60 以下、60--70、70--80、80--90、90 以上分组做频数表；接着画一张直方图看看成绩大致集中在哪一段；最后算出平均分、标准差，心里就有数了。

\begin{remark}
整理和描述不是“为了好看而折腾”。同一个数据集，分组方式、图表类型选得不一样，给人留下的印象可能完全不同。整理阶段多花点心思，能避免后面被数据“误导”。
\end{remark}



\section{计算与分析数据}
数据整理好以后，就到了真正“算”和“想”的环节。\textbf{计算与分析数据}，不只是把统计量按公式计算出来，更重要的是\textbf{结合研究问题去解读这些数字}：它们说明了什么？能不能支撑我们的猜想？

这一步通常按这样的思路走：
\begin{itemize}
  \item \textbf{选对统计量}：根据数据类型和研究问题，选择合适的手段。比如比较平均水平用均值或中位数，看离散程度用标准差，看两个变量的关系用相关系数。
  \item \textbf{计算描述统计量}：先用均值、中位数、标准差、分位数等把数据的“长相”描述出来。
  \item \textbf{进行统计推断}：如果需要由样本推测总体，就做点估计、区间估计或假设检验（如$t$检验、$F$检验、$\chi^2$检验）。
  \item \textbf{回到问题本身}：把计算结果放回最初的研究问题里，用自然语言下结论，而不是扔一堆数字让读者自己猜。
\end{itemize}

例如：你算出 50 名日语学生的平均分是 82 分、标准差是 8 分。光说“平均 82 分”还不够，你还想知道这个样本能否代表全校——于是做一次$t$检验，判断全校日语学生的平均分是否显著高于 80 分，再根据$p$值给出结论。

\begin{remark}
计算只是手段，不是终点。拿到一个“显著”的结果，也别急着欢呼：还要想想样本有没有代表性、效应在实际中够不够大、有没有其他解释。\textbf{统计显著不等于实际重要，更不等于因果关系。}
\end{remark}



\section{完整的统计步骤}
把前面几节串起来，一次完整的统计大致会经历下面几步。值得注意的是，这条路并不是只能走一遍：当分析结果不理想、发现设计有漏洞时，常常需要回到前面的步骤重新调整。

\begin{figure}[htbp]
  \centering
  \begin{tikzpicture}[
      node distance=1.1cm,
      box/.style={draw=main, very thick, rounded corners=8pt, fill=main!8,
                  minimum width=3.4cm, minimum height=0.9cm, align=center, font=\small},
      arrow/.style={-{Stealth[length=3mm]}, thick, main}
    ]
    \node[box] (q) {提出问题};
    \node[box, below=of q] (d) {统计设计};
    \node[box, below=of d] (s) {样本数据采集};
    \node[box, below=of s] (o) {数据的整理与描述};
    \node[box, below=of o] (a) {计算与分析数据};
    \node[box, below=of a] (c) {得出结论};

    \draw[arrow] (q) -- (d);
    \draw[arrow] (d) -- (s);
    \draw[arrow] (s) -- (o);
    \draw[arrow] (o) -- (a);
    \draw[arrow] (a) -- (c);

    \draw[arrow] (a.west) -- ++(-2.2,0) node[midway, left, font=\scriptsize, align=center]
          {结果不理想时\\回头调整} |- (q.west);
  \end{tikzpicture}
  \caption{完整的统计步骤流程图}
\end{figure}

\begin{remark}
流程图里最容易被忽略的是那条“回头”的箭头。现实中很少有统计能一次走通：数据收得不够、分析结果不显著、问题本身要修正，都是常有的事。\textbf{允许回头，才是完整的统计过程。}
\end{remark}









\chapter{数据的可视化与描述}

% \begin{introduction}
%   \item 积分定义~\ref{def:int}
%   \item Fubini 定理~\ref{thm:fubi}
%   \item 最优性原理~\ref{pro:max}
%   \item 柯西列性质~\ref{property:cauchy}
%   \item 韦达定理
% \end{introduction}

\section{使用图展示数据}
在前面的章节做了必要的准备后，本章开始介绍。数据整理好以后，一张合适的图往往比一大段文字更能让人一眼看出规律。本章先用\textbf{二维图}入门，再逐步介绍三维图和降维显示。

\subsection{二维图}
二维图是在一张平面图上，用横轴和纵轴承载信息，最常见的有直方图、折线图、箱型图和散点图。

\subsubsection{直方图}
直方图用来展示\textbf{一组数值数据的分布}：把数据按区间分组，每个区间画一根柱子，柱子的高度表示落在该区间的数据个数（频数）。

例如，某次日语测验 40 名学生的成绩如下（单位：分）：
\begin{center}
56, 58, 61, 63, 65, 67, 68, 70, 71, 72,\\
73, 74, 75, 75, 76, 77, 78, 78, 79, 80,\\
80, 81, 82, 83, 84, 85, 85, 86, 87, 88,\\
89, 90, 91, 92, 93, 95, 96, 97, 98, 100
\end{center}

按每 10 分一组，可以得到下面的频数表：
\begin{table}[htbp]
  \centering
  \caption{40 名学生的日语成绩频数表}
  \begin{tabular}{cc}
    \toprule
    成绩区间（分） & 人数 \\
    \midrule
    50--59 & 2 \\
    60--69 & 5 \\
    70--79 & 12 \\
    80--89 & 12 \\
    90--100 & 9 \\
    \bottomrule
  \end{tabular}
\end{table}

\begin{figure}[H]
  \centering
  \begin{tikzpicture}
    \begin{axis}[
        ybar,
        bar width=0.7,
        xlabel={成绩区间（分）},
        ylabel={人数},
        ymin=0,
        xtick={1,2,3,4,5},
        xticklabels={50--59,60--69,70--79,80--89,90--100},
        nodes near coords,
        every node near coord/.append style={font=\scriptsize},
        title={40 名学生日语成绩直方图}
      ]
      \addplot coordinates {(1,2) (2,5) (3,12) (4,12) (5,9)};
    \end{axis}
  \end{tikzpicture}
  \caption{40 名学生日语成绩直方图}
\end{figure}

从直方图可以直观看到：成绩大多集中在 70--89 分，两头少、中间多，整体分布比较接近“中间高、两边低”的形状。

\subsubsection{折线图}
折线图适合展示\textbf{数据随时间或其他有序变量的变化趋势}。横轴通常是有先后顺序的类别（如星期、月份），纵轴是数值。

例如，某位同学连续一周每天复习日语的时间（单位：分钟）如下：
\begin{center}
周一 30，周二 45，周三 40，周四 60，周五 50，周六 90，周日 80
\end{center}

\begin{figure}[htbp]
  \centering
  \begin{tikzpicture}
    \begin{axis}[
        symbolic x coords={周一,周二,周三,周四,周五,周六,周日},
        xtick=data,
        xlabel={星期},
        ylabel={复习时间（分钟）},
        ymin=0,
        mark=*,
        title={一周每天复习日语的时间}
      ]
      \addplot coordinates {(周一,30) (周二,45) (周三,40) (周四,60) (周五,50) (周六,90) (周日,80)};
    \end{axis}
  \end{tikzpicture}
  \caption{一周每天复习日语时间折线图}
\end{figure}

从折线图可以清楚看到：周中起伏不大，但到了周末复习时间明显上升——这就是折线图“看趋势”的用处。

\subsubsection{箱型图}
箱型图用来\textbf{概括一组数据的分布位置和离散程度}，它能同时显示最小值、下四分位数、中位数、上四分位数和最大值，还能把异常值“暴露”出来。

例如，比较三个班的日语成绩（每班 10 人）：
\begin{center}
1班：60, 65, 70, 72, 75, 78, 80, 82, 85, 90\\
2班：55, 60, 62, 68, 70, 72, 75, 78, 80, 85\\
3班：70, 75, 78, 82, 85, 88, 90, 92, 95, 98
\end{center}

\begin{figure}[H]
  \centering
  \begin{tikzpicture}
    \begin{axis}[
        ylabel={成绩（分）},
        ymin=40, ymax=110,
        xtick={1,2,3},
        xticklabels={1班,2班,3班},
        boxplot/draw direction=y,
        title={三个班日语成绩箱型图}
      ]
      \addplot+[boxplot prepared={
          median=76.5, upper quartile=82, lower quartile=70,
          upper whisker=90, lower whisker=60
        }, boxplot/draw position=1] coordinates {};
      \addplot+[boxplot prepared={
          median=71, upper quartile=78, lower quartile=62,
          upper whisker=85, lower whisker=55
        }, boxplot/draw position=2] coordinates {};
      \addplot+[boxplot prepared={
          median=86.5, upper quartile=92, lower quartile=78,
          upper whisker=98, lower whisker=70
        }, boxplot/draw position=3] coordinates {};
    \end{axis}
  \end{tikzpicture}
  \caption{三个班日语成绩箱型图}
\end{figure}

一眼看去，3 班的箱子整体位置最高，说明成绩整体最好；2 班的箱子最“胖”，说明成绩比较分散。

\subsubsection{散点图}
散点图用来\textbf{观察两个数值变量之间是否有关系}，每个点代表一条数据。

例如，10 名同学每天学习日语的时间与日语成绩如下：
\begin{center}
学习时间（小时）：2, 3, 4, 5, 6, 7, 8, 9, 10, 11\\
日语成绩（分）：58, 62, 68, 66, 75, 80, 83, 96, 92, 95
\end{center}

\begin{figure}[htbp]
  \centering
  \begin{tikzpicture}
    \begin{axis}[
        xlabel={每天学习时间（小时）},
        ylabel={日语成绩（分）},
        xmin=0, xmax=12,
        ymin=50, ymax=100,
        only marks,
        title={学习时间与日语成绩的散点图}
      ]
      \addplot coordinates {(2,58) (3,62) (4,68) (5,66) (6,75) (7,80) (8,83) (9,96) (10,92) (11,95)};
    \end{axis}
  \end{tikzpicture}
  \caption{学习时间与日语成绩散点图}
\end{figure}

这些点大致从左下往右上排列，说明学习时间越长、成绩越高，两者呈现正相关——更严格的度量方法（相关系数）会在后面章节介绍。

\begin{remark}
二维图虽然简单，但选错图很容易误导读者：想看分布用直方图，想看趋势用折线图，想比分布位置用箱型图，想看两个变量的关系用散点图。\textbf{先想清楚要回答什么问题，再选图进行展示。}
\end{remark}

\subsection{三维图}
二维图只能同时看两个变量，当我们想观察\textbf{三个变量}的关系时，就可以把图“立起来”，变成三维图。三维图最常见的形式是三维曲面图和三维散点图，其中\textbf{三维曲面图}特别适合展示“一个数值随两个变量连续变化”的情况。

例如，二元正态分布的密度曲面就像一个“山包”：中心最高，向四周平滑地下降。这张图用 $z=e^{-(x^2+y^2)/2}$ 示意，横轴和纵轴是两个变量，竖轴是密度值。

\begin{figure}[H]
  \centering
  \begin{tikzpicture}
    \begin{axis}[
        view={60}{30},
        xlabel={$x$},
        ylabel={$y$},
        zlabel={密度值},
        zmin=0, zmax=1.1,
        title={二元正态密度曲面}
      ]
      \addplot3[surf, domain=-3:3, domain y=-3:3, samples=30]
        {exp(-(x^2+y^2)/2)};
    \end{axis}
  \end{tikzpicture}
  \caption{二元正态密度曲面示意图}
\end{figure}

三维曲面图能让人一眼看到“整体形状”，比如哪里高、哪里低、怎么过渡。但它也有代价：在纸面上不容易精确读数，人眼也容易看花。所以实际分析中，三维图多用来“找感觉”，精确判断还是要靠统计量。

\subsection{高维度的降维显示}
变量一旦超过三个，二维、三维图都画不下了。这时常用的思路是\textbf{降维}：把高维数据投影到二维或三维空间，同时尽量保留数据原有的差异和结构。

常见的降维方法有：
\begin{itemize}
  \item \textbf{主成分分析（PCA）}：找几个新的、互不相关的“综合变量”，用它们代表原来的众多变量。
  \item \textbf{t-SNE}：擅长把高维数据“摊开”到二维，让相似的样本在图上靠得更近，常用于可视化聚类结果。
  \item \textbf{UMAP}：类似 t-SNE，但速度更快、对全局结构保留得更好。
\end{itemize}

\begin{remark}
降维后的图只是“投影”，不是原始数据的完整还原。它适合用来发现大致的模式（比如哪些样本聚在一起），但下结论之前，还是要回到原始变量上验证。\textbf{降维是辅助眼睛，不是替代统计。}
\end{remark}



\section{使用统计量描述数据}
图能让人“看见”数据，但写报告、做比较时，还是需要数字来精确地描述。这一节把第 2 章提到的统计量正式用起来，分成“常用统计量”和“各个分布的统计量”两部分。

\subsection{常用统计量}
常用统计量大致可以分成三类：描述\textbf{集中趋势}、描述\textbf{离散程度}、描述\textbf{分布位置与形状}。

\textbf{1. 集中趋势：数据“集中在哪”}
\begin{itemize}
  \item \textbf{样本均值} $\bar{x}=\dfrac{1}{n}\sum\limits_{i=1}^{n}x_i$：所有数据的算术平均，最常用，但容易被极端值拉偏。
  \item \textbf{中位数}：排序后正中间的数，极端值再多也“不为所动”。
  \item \textbf{众数}：出现次数最多的取值，适合描述分类数据或“最常见的水平”。
\end{itemize}

\textbf{2. 离散程度：数据“散不散”}
\begin{itemize}
  \item \textbf{极差} $R=x_{(n)}-x_{(1)}$：最大值减最小值，简单，但只看了两个极端值。
  \item \textbf{样本方差} $s^2=\dfrac{1}{n-1}\sum\limits_{i=1}^{n}(x_i-\bar{x})^2$：数据偏离均值的平均平方程度。
  \item \textbf{样本标准差} $s=\sqrt{s^2}$：和原始数据同单位，比方差更直观。
  \item \textbf{变异系数} $CV=\dfrac{s}{\bar{x}}$：去掉单位的影响，用来比较尺度不同的数据谁更“散”。
\end{itemize}

\textbf{3. 分布位置与形状}
\begin{itemize}
  \item \textbf{分位数}：如第 25\%、50\%、75\% 分位数（下四分位数、中位数、上四分位数），描述数据分布在哪些位置。
  \item \textbf{偏度}：分布是否左右对称、偏向哪边。
  \item \textbf{峰度}：分布是“尖”还是“平”、尾部厚不厚。
\end{itemize}

用回 4.1.1 那 40 名学生的日语成绩，可以算出下面这些常用统计量：
\begin{table}[H]
  \centering
  \caption{40 名学生成绩的常用统计量}
  \begin{tabular}{lc|lc}
    \toprule
    统计量 & 数值 & 统计量 & 数值 \\
    \midrule
    均值 $\bar{x}$ & 79.95 & 极差 $R$ & 44 \\
    中位数 & 80 & 样本方差 $s^2$ & 128.20 \\
    下四分位数 & 72.5 & 样本标准差 $s$ & 11.32 \\
    上四分位数 & 88.5 & 变异系数 $CV$ & 0.14 \\
    \bottomrule
  \end{tabular}
\end{table}

\begin{remark}
这份数据里，均值 79.95 和中位数 80 非常接近，说明成绩分布比较对称，没有明显的极端值捣乱；标准差约 11.32 分，大致可以理解为“大部分学生的成绩在平均分上下 10 分出头浮动”。\textbf{只看均值不看离散程度，很容易被“平均”骗了。}
\end{remark}

\subsection{各个分布的统计量}
不同的概率分布都有自己固定的期望（均值）和方差。记住这些结论，后面做参数估计和假设检验时能省很多事。

\begin{table}[H]
  \centering
  \caption{常见分布的期望与方差}
  \begin{tabular}{lccc}
    \toprule
    分布 & 记号 & 期望 $E(X)$ & 方差 $Var(X)$ \\
    \midrule
    0-1 分布 & $B(1,p)$ & $p$ & $p(1-p)$ \\
    二项分布 & $B(n,p)$ & $np$ & $np(1-p)$ \\
    泊松分布 & $P(\lambda)$ & $\lambda$ & $\lambda$ \\
    均匀分布 & $U(a,b)$ & $\frac{a+b}{2}$ & $\frac{(b-a)^2}{12}$ \\
    正态分布 & $N(\mu,\sigma^2)$ & $\mu$ & $\sigma^2$ \\
    指数分布 & $Exp(\lambda)$ & $\frac{1}{\lambda}$ & $\frac{1}{\lambda^2}$ \\
    $t$ 分布 & $t(k)$ & $0\ (k>1)$ & $\frac{k}{k-2}\ (k>2)$ \\
    $\chi^2$ 分布 & $\chi^2(k)$ & $k$ & $2k$ \\
    $F$ 分布 & $F(k_1,k_2)$ & $\frac{k_2}{k_2-2}\ (k_2>2)$ & $\frac{2k_2^2(k_1+k_2-2)}{k_1(k_2-2)^2(k_2-4)}\ (k_2>4)$ \\
    \bottomrule
  \end{tabular}
\end{table}

\begin{remark}
几个容易记的小规律：泊松分布“期望等于方差”；正态分布完全由 $\mu$ 和 $\sigma^2$ 两个参数决定；$t$、$\chi^2$、$F$ 这几个用于检验的分布，其期望和方差往往只在自由度足够大时才存在或才简洁。
\end{remark}






\chapter{概率与分布}
\section{概率为0就是不可能事件吗?}
直觉上，概率是 0 的事情好像就“绝不会发生”。但在统计学里，这个直觉要小心：\textbf{概率为 0 并不等于不可能发生}，\textbf{概率为 1 也不等于必然发生}。

关键在于区分\textbf{离散型}和\textbf{连续型}随机变量：
\begin{itemize}
  \item \textbf{离散型}：取值是一个个孤立的点，概率 0 通常就意味着“这件事不发生”。
  \item \textbf{连续型}：取值是某个区间里的任意实数，任何一个“单点”的概率都是 0。例如 $X\sim U(0,1)$，$P(X=0.5)=0$，但 0.5 并不是不可能出现——只是它太“细”了，单独一个点没有“长度”，所以概率算出来是 0。
\end{itemize}

\begin{remark}
可以这样理解：在数轴上随机撒一个点，它正好落在 0.5 这个“没有宽度”的位置上，概率当然是 0；但这并不意味着 0.5 不存在。所以连续分布里，“概率为 0 的事件”和“不可能事件”是两回事。
\end{remark}

\section{什么是分布?}
“分布”这个词，本质上是在回答一个问题：\textbf{一个随机变量取各种可能值的概率，是怎么“摊开”的}。

例如掷一颗公平的骰子：
\begin{itemize}
  \item 可能取值是 1、2、3、4、5、6；
  \item 每个点数出现的概率都是 $\frac{1}{6}$。
\end{itemize}
把“每个取值对应多少概率”完整画出来，就是下面这张柱状图：

\begin{figure}[H]
  \centering
  \begin{tikzpicture}
    \begin{axis}[
        ybar,
        bar width=0.5,
        ymin=0, ymax=0.35,
        xtick={1,2,3,4,5,6},
        ylabel={概率},
        xlabel={点数},
        nodes near coords,
        every node near coord/.append style={font=\scriptsize},
        title={公平骰子的概率分布}
      ]
      \addplot[point meta=explicit symbolic] coordinates {
        (1,1/6) [$\frac{1}{6}$]
        (2,1/6) [$\frac{1}{6}$]
        (3,1/6) [$\frac{1}{6}$]
        (4,1/6) [$\frac{1}{6}$]
        (5,1/6) [$\frac{1}{6}$]
        (6,1/6) [$\frac{1}{6}$]
      };
    \end{axis}
  \end{tikzpicture}
  \caption{公平骰子的概率分布（离散均匀分布）}
\end{figure}

这种取值是孤立点、每个点都有明确概率的，叫\textbf{离散型分布}。注意，所有柱子的高度加起来正好等于 1。

但现实中的骰子不一定是公平的。假如这颗骰子被动了手脚，1 点和 2 点更容易出现：
\begin{itemize}
  \item $P(1)=\frac{1}{4},\ P(2)=\frac{1}{4},\ P(3)=\frac{3}{20},\ P(4)=\frac{3}{20},\ P(5)=\frac{1}{10},\ P(6)=\frac{1}{10}$
\end{itemize}
画出来就是这样一张“高低不平”的图：

\begin{figure}[H]
  \centering
  \begin{tikzpicture}
    \begin{axis}[
        ybar,
        bar width=0.5,
        ymin=0, ymax=0.3,
        xtick={1,2,3,4,5,6},
        ylabel={概率},
        xlabel={点数},
        nodes near coords,
        every node near coord/.append style={font=\scriptsize},
        title={被动手脚后的骰子概率分布}
      ]
      \addplot[point meta=explicit symbolic] coordinates {
        (1,1/4) [$\frac{1}{4}$]
        (2,1/4) [$\frac{1}{4}$]
        (3,3/20) [$\frac{3}{20}$]
        (4,3/20) [$\frac{3}{20}$]
        (5,1/10) [$\frac{1}{10}$]
        (6,1/10) [$\frac{1}{10}$]
      };
    \end{axis}
  \end{tikzpicture}
  \caption{被动手脚后的骰子概率分布（不均匀分布）}
\end{figure}

同样是 1--6 的取值，概率怎么分配不一样，分布就不一样——这就是“分布”的意义。

所以一句话总结：\textbf{分布 = 把所有可能取值，以及它们“被抽中的可能性”，完整地描述出来}。

\begin{remark}
分布就像是数据的“户口本”：它告诉我们数据可能长什么样、哪些值常见、哪些值罕见。后面要讲的标准正态分布、二项分布、泊松分布等，都是具有某种特定“长相”的分布。
\end{remark}

\section{标准正态分布}
在讲标准正态分布之前，我们先要知道\textbf{连续型分布}是怎么回事。
与\textbf{取值}是孤立点、每个点都有明确概率的\textbf{离散型分布}不同，
我们有时会遇见\textbf{连续}的数据，这里的连续不是排队那种位置上的前后连续，
而是分布中，\textbf{数值取值}的连续，
比如我们现在有一个无数个面的概念骰子，
它能摇出1点，还能摇出1.999点,甚至是1.0000...1至6点之间所有的数。
像这样取值是连续的，可以无限接近的数据的分布，就是连续型分布。
正态分布是统计里最重要的连续型分布之一，而\textbf{标准正态分布}是其中最“标准”的一款：均值 $\mu=0$、方差 $\sigma^2=1$，记作 $N(0,1)$。它的概率密度函数为
\begin{center}
  $f(x)=\dfrac{1}{\sqrt{2\pi}}e^{-x^2/2}$
\end{center}

标准正态分布曲线长得像一个左右对称的“钟”：曲线下方的面积之和，即总概率"1"
\begin{figure}[H]
  \centering
  \begin{tikzpicture}
    \begin{axis}[
        width=10cm, height=6cm,
        xlabel={$x$},
        ylabel={密度值},
        ymin=0, ymax=0.45,
        xmin=-4.2, xmax=4.2,
        title={标准正态分布 $N(0,1)$ 密度曲线}
      ]
      \addplot[domain=-4:4, samples=200, thick, main]
        {1/sqrt(2*pi)*exp(-x^2/2)};
      \addplot[mark=none, dashed] coordinates {(-1,0) (-1,0.35)};
      \addplot[mark=none, dashed] coordinates {(1,0) (1,0.35)};
      \addplot[mark=none, dashed] coordinates {(-2,0) (-2,0.2)};
      \addplot[mark=none, dashed] coordinates {(2,0) (2,0.2)};
      \node at (axis cs:0,-0.06) {$0$};
      \node at (axis cs:-1,-0.06) {$-1$};
      \node at (axis cs:1,-0.06) {$1$};
      \node at (axis cs:-2,-0.06) {$-2$};
      \node at (axis cs:2,-0.06) {$2$};
    \end{axis}
  \end{tikzpicture}
  \caption{标准正态分布密度曲线}
\end{figure}

例如，假设某校学生的身高近似服从正态分布，均值 173 cm、标准差 5 cm。它的密度曲线仍然是一个“钟”，但对称轴从 0 挪到了 173 cm：

\begin{figure}[H]
  \centering
  \begin{tikzpicture}
    \begin{axis}[
        width=10cm, height=6cm,
        xlabel={身高（cm）},
        ylabel={密度值},
        ymin=0, ymax=0.09,
        xmin=155, xmax=191,
        title={身高 $N(173,5^2)$ 的正态分布曲线}
      ]
      \addplot[domain=155:191, samples=200, thick, main]
        {1/(5*sqrt(2*pi))*exp(-(x-173)^2/(2*25))};
      \addplot[mark=none, dashed] coordinates {(173,0) (173,0.08)};
      \node at (axis cs:173,-0.004) {$173$};
    \end{axis}
  \end{tikzpicture}
  \caption{以身高为例的正态分布曲线（对称轴 173 cm）}
\end{figure}

分布曲线关于 173 cm 完全对称，越靠近 173 cm 密度越高，越往两边越低。
这表示这个学校绝大多数学生的身高都"中规中矩"，在173cm，
长得特别高的人数(或概率)比较小，同样的，长得特别矮的人数(或概率)也比较小。
换句话说，一组符合正态分布的数据，绝大部分都中规中矩，越是极端的数据，出现概率越小。

标准正态分布有几个常用性质：
\begin{itemize}
  \item 关于 $x=0$ 完全对称，均值、中位数、众数都等于 0。
  \item 曲线下的总面积等于 1，也就是所有可能取值的概率总和为 1。
  \item \textbf{68--95--99.7 法则}：约 68\% 的数据落在正负一个标准差内，约 95\% 落在正负两个标准差内，约 99.7\% 落在正负三个标准差内。
\end{itemize}

\begin{remark}
任何正态分布 $N(\mu,\sigma^2)$ 都可以通过标准化
\begin{center}
  $Z=\dfrac{X-\mu}{\sigma}$
\end{center}
变成标准正态分布 $N(0,1)$。所以查概率表时，只需要一张标准正态分布表就够了。这也是标准正态分布如此“标准”的原因。
\end{remark}



\section{离散型随机变量及其分布}
离散型随机变量的取值是\textbf{一个接一个的孤立点}，每个取值都有明确的概率，所有取值的概率加起来等于 1。前面 5.2 里掷骰子就是典型的离散型分布。下面介绍三个最常用的离散分布：0-1分布、二项分布、泊松分布。

\subsection{0-1分布}
\textbf{0-1分布}是最简单的离散分布，随机变量只取两个值：0 和 1。通常把“成功”记为 1、“失败”记为 0，成功概率为 $p$：
\begin{center}
$P(X=1)=p,\qquad P(X=0)=1-p$
\end{center}

例如，抛一次硬币，正面朝上记为 1、反面记为 0，就服从 0-1 分布；一次考试及格记为 1、不及格记为 0，也服从 0-1 分布。

0-1 分布的期望和方差都很简洁：
\begin{itemize}
  \item 期望：$E(X)=p$
  \item 方差：$Var(X)=p(1-p)$
\end{itemize}

\subsection{二项分布}
把 0-1 试验独立地重复 $n$ 次，记录“成功”的总次数 $X$，$X$ 就服从\textbf{二项分布}，记作 $X\sim B(n,p)$。它刻画的是：在 $n$ 次独立重复试验中恰好成功 $k$ 次的概率
\begin{center}
$P(X=k)=\binom{n}{k}p^k(1-p)^{n-k},\qquad k=0,1,\ldots,n$
\end{center}

例如，连续抛 10 次硬币，正面朝上的次数 $X\sim B(10,0.5)$。正面出现 0 次、5 次、10 次的概率各不相同，画出来就是下面这张柱状图：

\begin{figure}[H]
  \centering
  \begin{tikzpicture}
    \begin{axis}[
        ybar,
        bar width=0.6,
        ymin=0, ymax=0.3,
        xtick={0,1,2,3,4,5,6,7,8,9,10},
        ylabel={概率},
        xlabel={正面朝上次数 $k$},
        title={二项分布 $B(10,0.5)$}
      ]
      \addplot coordinates {
        (0,0.00098) (1,0.00977) (2,0.04395) (3,0.11719) (4,0.20508) (5,0.24609)
        (6,0.20508) (7,0.11719) (8,0.04395) (9,0.00977) (10,0.00098)
      };
    \end{axis}
  \end{tikzpicture}
  \caption{二项分布 $B(10,0.5)$ 的概率分布}
\end{figure}

二项分布的期望和方差：
\begin{itemize}
  \item 期望：$E(X)=np$
  \item 方差：$Var(X)=np(1-p)$
\end{itemize}

\begin{remark}
0-1 分布其实就是 $n=1$ 时的二项分布。二项分布适合“次数有限、每次成功概率固定”的独立重复试验。
\end{remark}

\subsection{泊松分布}
\textbf{泊松分布}用来描述单位时间、单位面积或单位长度内“稀有事件”发生次数的分布，记作 $X\sim P(\lambda)$，其中 $\lambda$ 是平均发生次数：
\begin{center}
$P(X=k)=\dfrac{\lambda^k e^{-\lambda}}{k!},\qquad k=0,1,2,\ldots$
\end{center}

例如，一小时内到达某商店的顾客数、一页书上印刷错误的数量、某路段一小时内发生的交通事故数，都常用泊松分布描述。以 $\lambda=3$ 为例，画出来就是下面这张柱状图：

\begin{figure}[H]
  \centering
  \begin{tikzpicture}
    \begin{axis}[
        ybar,
        bar width=0.6,
        ymin=0, ymax=0.25,
        xtick={0,1,2,3,4,5,6,7,8},
        ylabel={概率},
        xlabel={事件发生次数 $k$},
        title={泊松分布 $P(3)$}
      ]
      \addplot coordinates {
        (0,0.0498) (1,0.1494) (2,0.2240) (3,0.2240) (4,0.1680) (5,0.1008) (6,0.0504) (7,0.0216) (8,0.0081)
      };
    \end{axis}
  \end{tikzpicture}
  \caption{泊松分布 $P(3)$ 的概率分布}
\end{figure}

泊松分布最显眼的特征是：
\begin{itemize}
  \item 期望：$E(X)=\lambda$
  \item 方差：$Var(X)=\lambda$
\end{itemize}
期望等于方差，是泊松分布的“招牌”。

\begin{remark}
当二项分布的 $n$ 很大、$p$ 很小时，二项分布可以近似看成泊松分布（此时 $\lambda=np$）。所以泊松分布也常用来近似“大量试验中成功次数很少”的情形。
\end{remark}



\section{自由度}
\textbf{自由度}（degrees of freedom，简写 df）是一个听起来很高大上，其实很朴素的概念。不妨先来思考一个小问题：

如果已知一个 6 人小组的平均成绩，那么至少需要知道几个人的成绩，才能准确说出每个人的成绩？

答案是 5 个。因为在知道前 5 个人的成绩之前，最后一个人的成绩都是“未知”的；可一旦前 5 个人的成绩都知道了，最后一个人的成绩就被“锁死”了——它必须负责“兜底”，保证平均分对得上，这样才能“万事大吉”。

也就是说，这 6 个人里，真正能自由变化的只有前 $n-1=5$ 个人，剩下的 1 个人是“工具人”，专门用来保证平均值成立。这里的 $n-1$（样本量减 1）就是\textbf{自由度}，记作 df。

推广到一般情况：
\begin{itemize}
  \item \textbf{自由度 = 数据个数 − 不得不“牺牲”掉的约束个数}；
  \item 每估计一个参数（比如均值），通常就要少掉一个自由度。
\end{itemize}

这也是为什么样本方差
\begin{center}
$s^2=\dfrac{1}{n-1}\sum_{i=1}^{n}(x_i-\bar{x})^2$
\end{center}
要用 $n-1$ 做分母，而不是 $n$：因为均值 $\bar{x}$ 已经被估计出来了，数据不再完全自由。

常见分布里的自由度：
\begin{itemize}
  \item $\chi^2(k)$ 分布有 $k$ 个自由度；
  \item $t(k)$ 分布的自由度为 $k$；
  \item $F(k_1,k_2)$ 分布有两个自由度：分子自由度和分母自由度；
  \item 卡方独立性检验中，$r\times c$ 列联表的自由度为 $(r-1)(c-1)$。
\end{itemize}

\begin{remark}
自由度不是玄学，它决定了 $t$、$\chi^2$、$F$ 这些分布的“长相”。样本量越小、约束越多，自由度就越小，分布也就越“胖”、越不自信；样本量越大，自由度越大，这些分布会越来越接近正态分布。
\end{remark}






\chapter{假设检验：从统计量到结论}
\section{假设检验与$p$值解读}
先来想一个有点“抬杠”的问题：如何证明“并非世界上所有的人都是女人”？

直接证明它很麻烦，但\textbf{证伪}就容易多了——只需找到一个反例：一个男人。于是我们先把“世界上所有的人都是女人”当作一个待检验的假设（原假设 $H_0$），然后努力找反例；只要找到一个男人，就有充分理由推翻 $H_0$，从而接受“并非所有人都是女人”这个结论。

假设检验就是这个思路：\textbf{先假设一个“无事发生”的立场，再拿数据看能不能把它推翻}。

一套完整的假设检验通常按下面几步走：
\begin{enumerate}
  \item \textbf{提出假设}：原假设 $H_0$（通常是“没有差异/没有关系/没有效果”），备择假设 $H_1$（通常是“有差异/有关系/有效果”）。
  \item \textbf{选择合适的检验}：根据数据类型和问题选 $t$ 检验、$\chi^2$ 检验、$F$ 检验等。
  \item \textbf{计算检验统计量和 $p$ 值}：$p$ 值表示“在 $H_0$ 为真的前提下，出现当前样本结果（或更极端结果）的概率”。
  \item \textbf{比较显著性水平}：临界(概率)值通常取 $\alpha=0.05$。若 $p<\alpha$，说明在 $H_0$ 成立时还能看到这么极端的结果概率很小，于是有理由拒绝 $H_0$。
  \item \textbf{下结论}：如果计算出来 $p$ 值小于临界值，则拒绝 $H_0$ ，支持 $H_1$；若 $p$ 值高于临界值，不能直接接受 $H_0$ 只能说“证据不足”，不能说明原假设 $H_0$ 一定正确。
\end{enumerate}

举个例子：某张期末成绩表里，乍一看男同学的体育成绩比女同学平均高5分，现在想探究男女体育成绩是否真的存在差异。我们需要给出假设：
\begin{itemize}
  \item $H_0$：体育成绩与性别无关（差异只是巧合）；
  \item $H_1$：体育成绩与性别有关（存在真实差异）。
\end{itemize}
用统计软件或者手算在$H_0$前提下，得到现有成绩表的概率 $p$ 值：
\begin{itemize}
  \item 若 $p<0.05$：说明在“无关”的假设下，出现“男比女平均高 5 分”这种结果的可能性不到 5\%，我们更有理由相信这不是巧合，于是拒绝 $H_0$，认为体育成绩与性别存在差异。
  \item 若 $p=0.33$：说明在“无关”的假设下，出现这种结果的概率高达 33\%，这说明一切很可能只是巧合，于是不能拒绝 $H_0$。
\end{itemize}

\begin{remark}
$p$ 值小，说明“如果 $H_0$ 是真的，那这次样本结果太离谱了”，所以我们倾向于不相信 $H_0$。但要注意：\textbf{$p$ 值大不代表 $H_0$ 正确}，只代表“证据不足以拒绝它”。另外，$p<0.05$ 也不等于“实际效果很大”，更不等于因果关系。
\end{remark}



\section{置信区间}
用样本估计总体参数，最朴素的办法是给一个“点”：比如抽 100 个学生，算出平均身高 168 cm，就猜全校平均身高是 168 cm。但这样做太“自信”了——换一批样本，结果可能就不是 168 了。

更稳妥的做法是给一个\textbf{区间}：例如“有 95\% 的把握，全校平均身高在 166.2 cm 到 169.8 cm 之间”。这个区间就叫\textbf{置信区间}，95\% 叫\textbf{置信水平}。

\begin{definition}[置信区间]
设总体参数为 $\theta$，由样本构造区间 $(\hat{\theta}_L,\hat{\theta}_U)$。若该区间以 $1-\alpha$ 的概率包含 $\theta$，则称它为 $\theta$ 的置信水平为 $1-\alpha$ 的置信区间。
\end{definition}

常用的形式是“点估计 ± 误差范围”：
\begin{center}
$\bar{x}\pm z_{\alpha/2}\cdot\dfrac{\sigma}{\sqrt{n}}$（总体方差已知时）\qquad
$\bar{x}\pm t_{\alpha/2}(n-1)\cdot\dfrac{s}{\sqrt{n}}$（总体方差未知时）
\end{center}

怎么理解“95\%”？
\begin{itemize}
  \item 不是“总体均值有 95\% 的概率落在这个区间里”（总体参数是固定值，不是随机变量）；
  \item 而是：如果重复抽样很多次，每次都构造一个 95\% 置信区间，那么大约有 95\% 的区间会包含总体均值。
\end{itemize}

例如：抽 100 名学生，样本均值 168 cm，标准差 8 cm，则 95\% 置信区间约为
\begin{center}
$168\pm 1.96\times\dfrac{8}{\sqrt{100}}=168\pm1.568=(166.43,\ 169.57)$
\end{center}
可以写成：全校学生平均身高的 95\% 置信区间为 $(166.43,\ 169.57)$ cm。

\begin{remark}
置信区间比单个点估计更有信息量：区间越窄，说明估计越精确；区间越宽，说明不确定性越大。它和假设检验其实是一体两面——若某个假设值不在置信区间内，通常对应的检验也会拒绝 $H_0$。
\end{remark}

\section{$t$检验：关注“平均值”的差异}
\textbf{$t$ 检验}用来回答“平均值有没有差异”的问题，常见有三种：
\begin{itemize}
  \item \textbf{单样本 $t$ 检验}：检验一个样本的均值是否等于某个已知值。例如：50 名日语学生的平均分是否显著高于 80 分。
  \item \textbf{独立样本 $t$ 检验}：比较两个独立组的均值。例如：男、女生的体育成绩是否有差异。
  \item \textbf{配对样本 $t$ 检验}：比较同一批对象前后两次测量。例如：同一批学生培训前后的日语成绩是否有提高。
\end{itemize}

单样本 $t$ 检验的统计量为
\begin{center}
$t=\dfrac{\bar{x}-\mu_0}{s/\sqrt{n}}$
\end{center}
其中 $\bar{x}$ 是样本均值，$\mu_0$ 是假设的总体均值，$s$ 是样本标准差，$n$ 是样本量，$t$ 服从自由度为 $n-1$ 的 $t$ 分布。

例如：50 名学生日语成绩均值 82 分、标准差 8 分，想检验是否显著高于 80 分：
\begin{itemize}
  \item $H_0:\mu=80$，$H_1:\mu>80$；
  \item $t=\dfrac{82-80}{8/\sqrt{50}}\approx1.77$；
  \item 查 $t$ 分布表（df=49）得 $p$ 值，若 $p<0.05$ 则拒绝 $H_0$。
\end{itemize}

\begin{remark}
$t$ 检验适合小样本，且要求数据近似正态、两组方差不要差太多。样本量很大时，$t$ 分布会越来越接近标准正态分布，检验结果也会越来越“稳”。
\end{remark}

\section{$\chi^2$(卡方)检验：关注“数量/频率”的分布}
\textbf{$\chi^2$ 检验}（卡方检验）用于\textbf{分类数据}，关心的是“实际频数”和“理论频数”是否吻合。常见两种用途：
\begin{itemize}
  \item \textbf{拟合优度检验}：检验一组分类数据的分布是否符合某个理论分布。例如：一枚骰子是否公平。
  \item \textbf{独立性检验}：检验两个\textbf{分类变量}是否相互独立。例如：性别与是否选修日语课是否有关。
\end{itemize}

$\chi^2$ 统计量的基本形式：
\begin{center}
$\chi^2=\sum\dfrac{(O-E)^2}{E}$
\end{center}
其中 $O$ 是实际观测频数，$E$ 是理论期望频数。$O$ 和 $E$ 差得越多，$\chi^2$ 越大，越说明“不符合/不独立”。

例如检验一枚骰子是否公平：掷 60 次，理论上每个点数应出现 10 次；若 1 点出现 20 次、6 点出现 2 次，$O$ 和 $E$ 差很多，$\chi^2$ 就会很大，$p$ 值很小，从而怀疑骰子被动了手脚。

独立性检验的自由度：
\begin{center}
$df=(r-1)(c-1)$
\end{center}
其中 $r$ 是行数，$c$ 是列数。

\begin{remark}
$\chi^2$ 检验用的是“频数/次数”，不是“均值”。所以它和 $t$ 检验的分工很清楚：\textbf{比较平均值用 $t$ 检验，比较数量/频率的分布用 $\chi^2$ 检验}。
\end{remark}

\section{非技术用语的假设检验总结模板}
本节提供一个假设检验结果书写模板，便于简要地呈现假设检验过程与结论。使用时只需替换方括号中的内容。

\begin{quote}
本研究旨在探讨[研究目的]。据此提出如下假设：

\textbf{零假设（$H_0$）}：[没有差异/没有关系/没有效果]。
\textbf{备择假设（$H_1$）}：[存在差异/存在关系/存在效果]。

为检验上述假设，采用[检验方法，如独立样本 $t$ 检验、$\chi^2$ 独立性检验、单因素方差分析等]。检验的显著性水平设为 $\alpha=0.05$。

分析结果显示，检验统计量为[统计量值]，自由度为[自由度]，查表(或软件自动计算)知，对应的 $p$ 值为[$p$ 值]。

由于 $p$ 值[ ＜ / ＞ ]显著性水平 $\alpha=0.05$，[拒绝/不能拒绝]零假设 $H_0$。据此可以认为，[根据 $H_1$ 表述研究结论，如：样本数据提供了统计学上显著的证据，表明……；或：样本数据未提供足够的证据表明……]。

综上所述，[用一句话概括研究结论，但不作因果推断]。
\end{quote}

\begin{remark}
正式书写时有几点需要注意：
\begin{itemize}
  \item $p$ 值应给出具体数值或明确的比较关系（如 $p<0.05$、$p=0.032$），不宜只写“显著/不显著”；
  \item “拒绝 $H_0$”只能说“差异（或关系）具有统计学意义”，不能直接说“具有实际意义”或“存在因果关系”；
  \item 若不能拒绝 $H_0$，应表述为“没有充分的证据拒绝 $H_0$”，而不是“$H_0$ 成立”；
  \item 报告时应同时给出检验方法、统计量、自由度和 $p$ 值，保证结果可复核。
\end{itemize}
\end{remark}






\chapter{方差分析}
\section{前置知识--方差分析在“分析”什么}
$t$ 检验适合比较\textbf{两个}组的均值。那如果现在有四个班、两种教学方法、三种教材，要比较学生成绩的平均值呢？

一个天真的想法是两两做 $t$ 检验：4 个班之间要比较 6 次，3 种教材之间要比较 3 次。但比较次数一多，犯“假阳性”错误（本来没差异却说有差异）的概率就会悄悄变大，这叫做\textbf{多重比较问题}。

\textbf{方差分析}（ANOVA）就是为了解决这个问题而生的：它\textbf{一次性比较多个组的均值是否相等}，而不是两两乱比。

那么方差分析到底在“分析”什么？名字里有“方差”，但它关心的其实是\textbf{均值}。它的思路是把数据的总差异拆成两部分：
\begin{itemize}
  \item \textbf{组间差异}：各组均值之间的差异，反映“不同组是否真的不一样”；
  \item \textbf{组内差异}：同一组内部个体之间的差异，反映“随机波动有多大”。
\end{itemize}

如果组间差异相对组内差异足够大，说明“分组”确实起了作用，各组均值不像是一家人。这个“够不够大”用 $F$ 统计量来衡量：
\begin{center}
$F=\dfrac{\text{组间均方}}{\text{组内均方}}$
\end{center}
$F$ 越大，对应$p$ 值越小，越有理由认为各组均值存在显著差异。

\begin{remark}
方差分析有几个基本前提：各组数据近似正态、各组方差大致相等（方差齐性）、各组观测相互独立。前提满足得越好，结果越可靠。
\end{remark}

\section{单因素方差分析}
\textbf{单因素方差分析}只考察\textbf{一个因素}，比较这个因素的多个水平（组）之间均值是否有差异。

例如：想比较三种日语教材的教学效果，首先随机把学生分成 3 组，
每组用不同教材学习一学期，最后比较三组的平均成绩。
(如果学生已经按照不同教材学习了一学期，则直接收集数据即可，不用从零开始)
\begin{itemize}
  \item $H_0$：三组平均成绩相等（即三种教材的教学效果没有差异）；
  \item $H_1$：至少有一组的平均成绩与其他组不同（即至少一种教材的教学效果与其他教材有差异）。
\end{itemize}

方差分析的结果通常整理成下面这样一张方差分析表，也可由统计软件自动得出：

\begin{table}[H]
  \centering
  \caption{单因素方差分析表}
  \begin{tabular}{lcccc}
    \toprule
    变异来源 & 平方和 & 自由度 & 均方 & $F$ \\
    \midrule
    组间 & $SS_B$ & $k-1$ & $MS_B=\frac{SS_B}{k-1}$ & $\frac{MS_B}{MS_W}$ \\
    组内 & $SS_W$ & $n-k$ & $MS_W=\frac{SS_W}{n-k}$ & \\
    总计 & $SS_T$ & $n-1$ & & \\
    \bottomrule
  \end{tabular}
\end{table}

其中 $k$ 是组数，$n$ 是总样本量。软件会算出对应的 $p$ 值：若 $p<0.05$，拒绝 $H_0$，认为至少有一组均值不同。

\begin{remark}
方差分析只能告诉你“有没有差异”，不能告诉你“具体哪两组有差异”。想知道哪两组不同，还要再做\textbf{事后检验}（多重比较），比如 Tukey 检验。
\end{remark}

\section{双因素方差分析}
\textbf{双因素方差分析}同时考察\textbf{两个因素}对结果的影响，不仅可以检验每个因素的\textbf{主效应}，还能检验两个因素之间是否存在\textbf{交互效应}。

例如：研究“教学法”（因素 A，3 种）和“性别”（因素 B，2 种）对日语成绩的影响，可以得到 3$\times$2=6 个组合。双因素方差分析会分别检验：
\begin{itemize}
  \item 因素 A 的主效应：不同教学法之间成绩是否有差异；
  \item 因素 B 的主效应：男女之间成绩是否有差异；
  \item \textbf{交互效应}：教学法的效果是否“看人下菜”，比如某种教学法对女生特别有效、对男生却一般。
\end{itemize}

交互效应是双因素方差分析最值得关注的地方：如果两个因素之间存在交互作用，就不能只单独看某一个因素的“平均效果”，因为一个因素的效果可能依赖于另一个因素的水平。

\begin{remark}
画一张“交互图”很有帮助：横轴放一个因素，不同折线代表另一个因素的不同水平。若折线明显交叉或分岔，往往意味着存在交互效应。
\end{remark}

\section{多因素方差分析}
\textbf{多因素方差分析}把双因素推广到三个或更多因素，原理完全一样：检验每个因素的\textbf{主效应}，以及因素之间的各种\textbf{交互效应}。

例如：研究“教学法 $\times$ 性别 $\times$ 课时长度”三个因素对日语成绩的影响，就是一个三因素方差分析。

因素一多，事情会变得复杂：
\begin{itemize}
  \item 交互项的数量迅速增加（两两交互、三阶交互……）；
  \item 需要的样本量更大，否则很多组合“没人”；
  \item 结果解释也更烧脑：高阶交互往往很难用一句话讲清楚。
\end{itemize}

实际研究中，多因素方差分析通常交给统计软件完成，分析思路和单因素、双因素一脉相承：\textbf{先看交互，再看主效应，最后回到实际问题下结论}。

\begin{remark}
因素越多不等于越高级。如果样本量不够或交互项解释不清，强行塞很多因素反而会让结论变得不可靠。\textbf{好的设计是“够用就好”，不是“越多越好”。}
\end{remark}





\chapter{线性回归}
\section{相关不蕴含因果}
什么是相关性？什么是因果？如果“只要 A 发生，B 就发生”，那 A 和 B 就是因果关系；而相关只表示二者\textbf{存在某种关联}，至于关联有多强、是谁引起谁、有没有第三者在背后捣乱，都还说不准。

举个经典例子：入冬后，某商场短款羽绒服的销量逐渐上升，与此同时，同地区医院的骨折人数也在增加。把两条折线画在一起，趋势几乎同步，相关系数可能还挺高。难道卖羽绒服会导致骨折吗？这显然很荒谬。

其实真正的原因是\textbf{天气变冷}：气温越低，羽绒服卖得越好；同时路面结冰，摔骨折的人也变多。羽绒服销量和骨折人数只是“碰巧一起变”，它们之间并没有因果关系。

所以做统计时要时刻记住三件事：
\begin{itemize}
  \item \textbf{相关不意味着因果}：就算两个变量相关到“离谱”，也不能直接说一个导致另一个。
  \item \textbf{小心伪相关}：冰淇淋销量和溺亡人数、医院数量和奢侈品消费量，这类“纸面上相关、实际上无关”的例子很多，背后往往藏着被忽略的第三因素（比如气温、经济发展水平）。
  \item \textbf{不要预设结论}：带着“我想要的结果”去分析，容易只看到支持自己的证据。应该从一个公正的问题出发，用数据说话。
\end{itemize}

\begin{remark}
相关只能告诉你“它们总是一起变”，因果则需要实验设计、控制变量等更严格的手段才能判断。\textbf{统计上显著的相关，离“因果”还差得远。}
\end{remark}

\section{线性相关与决定系数$R^2$}
当两个数值变量大致呈直线关系时，我们用\textbf{相关系数} $r$ 来度量线性相关的方向和强弱：
\begin{center}
$r\in[-1,1]$
\end{center}
\begin{itemize}
  \item $r>0$：正相关，一个变大另一个也变大；
  \item $r<0$：负相关，一个变大另一个变小；
  \item $|r|$ 越接近 1，线性关系越强；越接近 0，线性关系越弱。
\end{itemize}

光看 $r$ 还不够，我们常常更关心：\textbf{$x$ 的变化能解释 $y$ 的多少变化？}这就用到\textbf{决定系数}：
\begin{center}
$R^2=r^2$
\end{center}
$R^2$ 的取值范围是 $[0,1]$，它表示\textbf{因变量 $y$ 的变异中有多大比例能被自变量 $x$ 的线性关系解释}。

例如，$r=0.8$ 时，$R^2=0.64$，意思是“学习时间”这个变量大约能解释“成绩”变异的 64\%，剩下 36\% 来自其他因素或随机波动。

\begin{remark}
$R^2$ 高说明线性拟合得好，但\textbf{不说明 $x$ 就是 $y$ 的原因}。哪怕 $R^2=0.99$，也可能只是伪相关。相关和回归描述的是“关系”，不是“因果”。
\end{remark}

\section{单变量线性回归}
\textbf{单变量线性回归}（也叫一元线性回归）用\textbf{一个}自变量 $x$ 去预测因变量 $y$，模型写成：
\begin{center}
$y=a+bx+\varepsilon$
\end{center}
其中 $a$ 是截距，$b$ 是斜率，$\varepsilon$ 是随机误差。拟合出来的一条直线叫\textbf{回归线}：
\begin{center}
$\hat{y}=a+bx$
\end{center}

回归线怎么找？最常用的方法是\textbf{最小二乘法}：让所有样本点到回归线的竖直距离（残差）的平方和最小。

例如，用“每天学习时间”预测“日语成绩”，10 名同学的数据可以拟合成
\begin{center}
$\hat{y}=50.25+4.15x$
\end{center}
意思是：学习时间每增加 1 小时，成绩平均提高约 4.15 分；学习 0 小时时，成绩的“基准”约为 50.25 分。

\begin{figure}[H]
  \centering
  \begin{tikzpicture}
    \begin{axis}[
        xlabel={每天学习时间（小时）},
        ylabel={日语成绩（分）},
        xmin=1, xmax=12,
        ymin=50, ymax=100,
        title={学习时间与日语成绩的回归线}
      ]
      \addplot[only marks, main] coordinates {(2,58) (3,62) (4,68) (5,71) (6,75) (7,80) (8,83) (9,88) (10,92) (11,95)};
      \addplot[domain=2:11, thick] {50.2545+4.14545*x};
    \end{axis}
  \end{tikzpicture}
  \caption{学习时间与日语成绩的散点图及回归线}
\end{figure}

\begin{remark}
回归线只能描述“平均而言”的关系，不代表每个个体都严格落在线上。而且用回归做预测时，最好只在数据范围(2-11h)内外推一点点，别拿它预测“每天学 18 小时”这种离谱情况。
\end{remark}



\section{双变量线性回归}
\textbf{双变量线性回归}用\textbf{两个}自变量 $x_1,x_2$ 预测 $y$：
\begin{center}
$y=a+b_1x_1+b_2x_2+\varepsilon$
\end{center}

例如，除了“每天学习时间”，再加入“每周复习次数”一起预测日语成绩：
\begin{center}
$\hat{y}=a+b_1\cdot\text{学习时间}+b_2\cdot\text{复习次数}$
\end{center}

双变量回归和单变量回归最大的不同在于：每个系数 $b_1,b_2$ 表示的是\textbf{在控制另一个变量不变的情况下}，该变量每增加一个单位时 $y$ 的平均变化。这就是“控制变量”的思想。

\begin{itemize}
  \item $b_1$：复习次数不变时，学习时间每多 1 小时，成绩平均变化多少；
  \item $b_2$：学习时间不变时，复习次数每多 1 次，成绩平均变化多少。
\end{itemize}

\begin{remark}
如果两个自变量彼此高度相关（比如“学习时间”和“复习次数”本来就一起变），就很难把它们的贡献分清楚，这叫\textbf{多重共线性}。变量之间越“独立”，系数的解释越干净。
\end{remark}

\section{多变量线性回归}
\textbf{多变量线性回归}把自变量推广到 $k$ 个：
\begin{center}
$y=a+b_1x_1+b_2x_2+\cdots+b_kx_k+\varepsilon$
\end{center}

它的思想和双变量完全一样：每个系数都是在“其他变量不变”的前提下，该变量的单独效应。只是变量一多，事情会变复杂：
\begin{itemize}
  \item 需要的样本量更大，否则系数估计不稳定；
  \item 变量之间可能互相纠缠（多重共线性）；
  \item 塞太多没用的变量容易\textbf{过拟合}——在已有数据上拟合得很好，但换一批数据就“失灵”。
\end{itemize}

实际建模时，通常先用相关系数和散点图观察变量关系，再逐步筛选变量、检查残差，最后用 $R^2$、调整 $R^2$ 和 $p$ 值等评估模型。

\begin{remark}
模型不是“变量越多越好”。\textbf{简单、可解释、经得起新数据检验的模型，往往比塞满变量的复杂模型更有用。}
\end{remark}








\nocite{*}
\printbibliography[heading=bibintoc, title=\ebibname]
\appendix

\chapter{各分布的概率临界值表}
本附录收录了 $t$ 分布、$\chi^2$ 分布和 $F$ 分布的常用临界值表，供查阅。

\section{$t$分布}
$t$ 分布临界值表给出的是：在给定自由度 df 和上尾概率 $\alpha$ 下，满足 $P(T>t_{\alpha})=\alpha$ 的临界值 $t_{\alpha}$。查表时先找自由度所在的行，再找对应的 $\alpha$ 列。

\begin{table}[H]
  \centering
  \caption{$t$ 分布临界值表（上尾概率 $\alpha$）}
  \small
  \begin{tabular}{c|ccccc}
    \toprule
    df & 0.10 & 0.05 & 0.025 & 0.01 & 0.005 \\
    \midrule
    1 & 3.078 & 6.314 & 12.706 & 31.821 & 63.657 \\
    2 & 1.886 & 2.920 & 4.303 & 6.965 & 9.925 \\
    3 & 1.638 & 2.353 & 3.182 & 4.541 & 5.841 \\
    4 & 1.533 & 2.132 & 2.776 & 3.747 & 4.604 \\
    5 & 1.476 & 2.015 & 2.571 & 3.365 & 4.032 \\
    6 & 1.440 & 1.943 & 2.447 & 3.143 & 3.707 \\
    7 & 1.415 & 1.895 & 2.365 & 2.998 & 3.499 \\
    8 & 1.397 & 1.860 & 2.306 & 2.896 & 3.355 \\
    9 & 1.383 & 1.833 & 2.262 & 2.821 & 3.250 \\
    10 & 1.372 & 1.812 & 2.228 & 2.764 & 3.169 \\
    11 & 1.363 & 1.796 & 2.201 & 2.718 & 3.106 \\
    12 & 1.356 & 1.782 & 2.179 & 2.681 & 3.055 \\
    13 & 1.350 & 1.771 & 2.160 & 2.650 & 3.012 \\
    14 & 1.345 & 1.761 & 2.145 & 2.624 & 2.977 \\
    15 & 1.341 & 1.753 & 2.131 & 2.602 & 2.947 \\
    16 & 1.337 & 1.746 & 2.120 & 2.583 & 2.921 \\
    17 & 1.333 & 1.740 & 2.110 & 2.567 & 2.898 \\
    18 & 1.330 & 1.734 & 2.101 & 2.552 & 2.878 \\
    19 & 1.328 & 1.729 & 2.093 & 2.539 & 2.861 \\
    20 & 1.325 & 1.725 & 2.086 & 2.528 & 2.845 \\
    21 & 1.323 & 1.721 & 2.080 & 2.518 & 2.831 \\
    22 & 1.321 & 1.717 & 2.074 & 2.508 & 2.819 \\
    23 & 1.319 & 1.714 & 2.069 & 2.500 & 2.807 \\
    24 & 1.318 & 1.711 & 2.064 & 2.492 & 2.797 \\
    25 & 1.316 & 1.708 & 2.060 & 2.485 & 2.787 \\
    26 & 1.315 & 1.706 & 2.056 & 2.479 & 2.779 \\
    27 & 1.314 & 1.703 & 2.052 & 2.473 & 2.771 \\
    28 & 1.313 & 1.701 & 2.048 & 2.467 & 2.763 \\
    29 & 1.311 & 1.699 & 2.045 & 2.462 & 2.756 \\
    30 & 1.310 & 1.697 & 2.042 & 2.457 & 2.750 \\
    40 & 1.303 & 1.684 & 2.021 & 2.423 & 2.704 \\
    60 & 1.296 & 1.671 & 2.000 & 2.390 & 2.660 \\
    120 & 1.289 & 1.658 & 1.980 & 2.358 & 2.617 \\
    $\infty$ & 1.282 & 1.645 & 1.960 & 2.326 & 2.576 \\
    \bottomrule
  \end{tabular}
\end{table}

\section{$\chi^2$分布}
$\chi^2$ 分布临界值表给出的是：在给定自由度 df 和上尾概率 $\alpha$ 下，满足 $P(\chi^2>\chi^2_{\alpha})=\alpha$ 的临界值。

\begin{table}[H]
  \centering
  \caption{$\chi^2$ 分布临界值表（上尾概率 $\alpha$）}
  \small
  \begin{tabular}{c|ccccc}
    \toprule
    df & 0.10 & 0.05 & 0.025 & 0.01 & 0.001 \\
    \midrule
    1 & 2.706 & 3.841 & 5.024 & 6.635 & 10.828 \\
    2 & 4.605 & 5.991 & 7.378 & 9.210 & 13.816 \\
    3 & 6.251 & 7.815 & 9.348 & 11.345 & 16.266 \\
    4 & 7.779 & 9.488 & 11.143 & 13.277 & 18.467 \\
    5 & 9.236 & 11.070 & 12.833 & 15.086 & 20.515 \\
    6 & 10.645 & 12.592 & 14.449 & 16.812 & 22.458 \\
    7 & 12.017 & 14.067 & 16.013 & 18.475 & 24.322 \\
    8 & 13.362 & 15.507 & 17.535 & 20.090 & 26.125 \\
    9 & 14.684 & 16.919 & 19.023 & 21.666 & 27.877 \\
    10 & 15.987 & 18.307 & 20.483 & 23.209 & 29.588 \\
    11 & 17.275 & 19.675 & 21.920 & 24.725 & 31.264 \\
    12 & 18.549 & 21.026 & 23.337 & 26.217 & 32.910 \\
    13 & 19.812 & 22.362 & 24.736 & 27.688 & 34.528 \\
    14 & 21.064 & 23.685 & 26.119 & 29.141 & 36.123 \\
    15 & 22.307 & 24.996 & 27.488 & 30.578 & 37.697 \\
    16 & 23.542 & 26.296 & 28.845 & 32.000 & 39.252 \\
    17 & 24.769 & 27.587 & 30.191 & 33.409 & 40.790 \\
    18 & 25.989 & 28.869 & 31.526 & 34.805 & 42.312 \\
    19 & 27.204 & 30.144 & 32.852 & 36.191 & 43.820 \\
    20 & 28.412 & 31.410 & 34.170 & 37.566 & 45.315 \\
    21 & 29.615 & 32.671 & 35.479 & 38.932 & 46.797 \\
    22 & 30.813 & 33.924 & 36.781 & 40.289 & 48.268 \\
    23 & 32.007 & 35.172 & 38.076 & 41.638 & 49.728 \\
    24 & 33.196 & 36.415 & 39.364 & 42.980 & 51.179 \\
    25 & 34.382 & 37.652 & 40.646 & 44.314 & 52.620 \\
    26 & 35.563 & 38.885 & 41.923 & 45.642 & 54.052 \\
    27 & 36.741 & 40.113 & 43.194 & 46.963 & 55.476 \\
    28 & 37.916 & 41.337 & 44.461 & 48.278 & 56.892 \\
    29 & 39.087 & 42.557 & 45.722 & 49.588 & 58.301 \\
    30 & 40.256 & 43.773 & 46.979 & 50.892 & 59.703 \\
    \bottomrule
  \end{tabular}
\end{table}



\newpage
\section{$F$分布}
$F$ 分布临界值表给出的是：在给定分子自由度 $\mathrm{df}_1$、分母自由度 $\mathrm{df}_2$ 和显著性水平 $\alpha$ 下，满足 $P(F>F_{\alpha})=\alpha$ 的临界值。下表为 $\alpha=0.05$ 的常用临界值。

\begin{table}[H]
  \centering
  \caption{$F$ 分布临界值表（$\alpha=0.05$）}
  \small
  \begin{tabular}{c|cccccccccccccc}
    \toprule
    $\mathrm{df}_2\backslash\mathrm{df}_1$ & 1 & 2 & 3 & 4 & 5 & 6 & 8 & 10 & 12 & 15 & 20 & 30 & 60 & $\infty$ \\
    \midrule
    1 & 161.45 & 199.50 & 215.71 & 224.58 & 230.16 & 233.99 & 238.88 & 241.88 & 243.91 & 245.95 & 248.01 & 250.10 & 252.20 & 254.32 \\
    2 & 18.51 & 19.00 & 19.16 & 19.25 & 19.30 & 19.33 & 19.37 & 19.40 & 19.41 & 19.43 & 19.45 & 19.46 & 19.48 & 19.50 \\
    3 & 10.13 & 9.55 & 9.28 & 9.12 & 9.01 & 8.94 & 8.85 & 8.79 & 8.74 & 8.70 & 8.66 & 8.62 & 8.57 & 8.53 \\
    4 & 7.71 & 6.94 & 6.59 & 6.39 & 6.26 & 6.16 & 6.04 & 5.96 & 5.91 & 5.86 & 5.80 & 5.75 & 5.69 & 5.63 \\
    5 & 6.61 & 5.79 & 5.41 & 5.19 & 5.05 & 4.95 & 4.82 & 4.74 & 4.68 & 4.62 & 4.56 & 4.50 & 4.43 & 4.36 \\
    6 & 5.99 & 5.14 & 4.76 & 4.53 & 4.39 & 4.28 & 4.15 & 4.06 & 4.00 & 3.94 & 3.87 & 3.81 & 3.74 & 3.67 \\
    7 & 5.59 & 4.74 & 4.35 & 4.12 & 3.97 & 3.87 & 3.73 & 3.64 & 3.57 & 3.51 & 3.44 & 3.38 & 3.30 & 3.23 \\
    8 & 5.32 & 4.46 & 4.07 & 3.84 & 3.69 & 3.58 & 3.44 & 3.35 & 3.28 & 3.22 & 3.15 & 3.08 & 3.01 & 2.93 \\
    9 & 5.12 & 4.26 & 3.86 & 3.63 & 3.48 & 3.37 & 3.23 & 3.14 & 3.07 & 3.01 & 2.94 & 2.86 & 2.79 & 2.71 \\
    10 & 4.96 & 4.10 & 3.71 & 3.48 & 3.33 & 3.22 & 3.07 & 2.98 & 2.91 & 2.84 & 2.77 & 2.70 & 2.63 & 2.54 \\
    12 & 4.75 & 3.89 & 3.49 & 3.26 & 3.11 & 3.00 & 2.85 & 2.75 & 2.69 & 2.62 & 2.54 & 2.47 & 2.40 & 2.30 \\
    15 & 4.54 & 3.68 & 3.29 & 3.06 & 2.90 & 2.79 & 2.64 & 2.54 & 2.48 & 2.40 & 2.33 & 2.25 & 2.18 & 2.07 \\
    20 & 4.35 & 3.49 & 3.10 & 2.87 & 2.71 & 2.60 & 2.45 & 2.35 & 2.28 & 2.20 & 2.12 & 2.04 & 1.95 & 1.84 \\
    30 & 4.17 & 3.32 & 2.92 & 2.69 & 2.53 & 2.42 & 2.27 & 2.16 & 2.09 & 2.01 & 1.93 & 1.84 & 1.75 & 1.62 \\
    60 & 4.00 & 3.15 & 2.76 & 2.53 & 2.37 & 2.25 & 2.10 & 2.00 & 1.92 & 1.84 & 1.75 & 1.65 & 1.53 & 1.39 \\
    $\infty$ & 3.84 & 3.00 & 2.60 & 2.37 & 2.21 & 2.10 & 1.94 & 1.83 & 1.75 & 1.67 & 1.57 & 1.46 & 1.32 & 1.00 \\
    \bottomrule
  \end{tabular}
\end{table}



\end{document}
